\studytypesubsection{Benchmarking LLMs for Software Engineering Tasks}
\label{sec:benchmarking-llms-for-software-engineering-tasks}

Benchmarking studies measure LLM performance on standardized SE tasks against reference outputs and shared metrics.

\studytypeparagraph{Description}

In a benchmark, the reference outputs serve as ground truth, and the metrics measure how well an LLM's outputs match them. Typical tasks include code generation, code summarization, code completion, and code repair~\cite{yang2025code}, but also natural language processing tasks such as anaphora resolution (i.e., the task of identifying the referring expression of a word or phrase occurring earlier in the text). %, interesting for subfields such a Requirements Engineering.
% GROUNDING yang2025code: "Modern LLMs have achieved remarkable capabilities across a wide range of code-related tasks, including code completion, translation, repair, and generation"
Metrics may include general metrics for text generation, such as \emph{ROUGE}, \emph{BLEU}, or \emph{METEOR}~\cite{DBLP:journals/tosem/HouZLYWLLLGW24}, or task-specific metrics, such as \emph{CodeBLEU} for code generation. Benchmarking requires high-quality reference datasets.
% GROUNDING DBLP:journals/tosem/HouZLYWLLLGW24: "Additionally, ROUGE/ROUGE-L [...], METEOR [...], EM (Exact Match) [...] are used in specific studies to evaluate the quality and accuracy of generated code or natural language code descriptions."

\studytypeparagraph{Examples}

In SE, benchmarking may include evaluating an LLM's ability to produce accurate and reliable outputs for a given input (usually a task description, possibly accompanied by data obtained from curated real-world projects or from synthetic SE-specific datasets). \emph{RepairBench}~\cite{DBLP:conf/llm4code/SilvaM25}, for example, contains 574 buggy Java methods and their corresponding fixed versions, which can be used to evaluate the performance of LLMs in code repair tasks. It uses the \emph{Plausible@1} metric (i.e., the probability that the first generated patch passes all test cases) and the \emph{AST Match@1} metric (i.e., the probability that the abstract syntax tree of the first generated patch matches the ground truth patch).
% GROUNDING DBLP:conf/llm4code/SilvaM25: "we define a maximum of 0.2 USD per evaluated bug, or approx. $115.1 for a total of 574 bugs"
\emph{SWE-Bench}~\cite{DBLP:conf/iclr/JimenezYWYPPN24} is a more generic benchmark that contains 2,294 SE Python tasks extracted from GitHub pull requests.
% GROUNDING DBLP:conf/iclr/JimenezYWYPPN24: "We therefore introduce SWE-bench, an evaluation framework including 2294 software engineering problems drawn from real GitHub issues and corresponding pull requests across 12 popular Python repositories."
To score the LLM's task performance, the benchmark validates whether the generated patch successfully compiles and calculates the percentage of passed test cases. Meanwhile, \emph{HumanEval}~\cite{DBLP:journals/corr/abs-2107-03374} is often used to assess code generation. 
% GROUNDING DBLP:journals/corr/abs-2107-03374: "On HumanEval, a new evaluation set we release to measure functional correctness for synthesizing programs from docstrings"
